\section{Facial Performance Transfer}

The third demo moves away from body motion and focuses instead on facial performance, using Unreal's Live Link workflow and MetaHuman Performance assets \cite{epic_live_link,epic_metahuman_performance} together with clips from the RAVDESS dataset \cite{livingstone_russo_ravdess_2018,ravdess_dataset} to generate facial animation and apply it to two MetaHumans: my own self-avatar and a female alternative character. Compared with the first two demos, this one is less about manually reshaping motion and more about judging how well an automated performance pipeline works in practice.

I selected clips representing several emotional states, including happiness, sadness, anger, fear, and surprise. Keeping the same source actor across the chosen clips made the comparison cleaner, because the more interesting variation then came from the target MetaHuman rather than from switching performers. I chose RAVDESS over MEAD mainly for practical reasons: the RAVDESS download I used was about 500 MB for a pack, whereas MEAD was closer to 18 GB on average, so RAVDESS was much easier to work with within the time and storage limits of the project.

\subsection{Import and Processing Workflow}

The source data came from the RAVDESS dataset \cite{livingstone_russo_ravdess_2018,ravdess_dataset}, where each filename encodes modality, vocal channel, emotion, intensity, statement, repetition, and actor identity, so for this demo I kept that structure as controlled as possible by using clips that all shared the audio-visual speech format, the same statement, the same actor, and the strong emotional intensity variant. The only thing I changed between clips was the emotion itself, which meant the comparison stayed controlled instead of mixing different performances together.

\begin{itemize}[leftmargin=*]
    \item \texttt{01-01-03-02-02-01-13.mp4} for happy,
    \item \texttt{01-01-04-02-02-01-13.mp4} for sad,
    \item \texttt{01-01-05-02-02-01-13.mp4} for angry,
    \item \texttt{01-01-06-02-02-01-13.mp4} for fearful,
    \item and \texttt{01-01-08-02-02-01-13.mp4} for surprised.
\end{itemize}

Using the same actor and the same spoken line mattered here, because it meant the differences in the final result could be read mainly as differences in emotion and target-character interpretation rather than noise caused by changing performer, delivery, or recording conditions.

One of the most useful findings in this demo was that the pipeline could be made much simpler by relying on Live Link Capture Manager \cite{epic_live_link} instead of overcomplicating the setup. Importing clips through this system automatically generated the associated capture data, media source, and audio assets, which reduced manual preparation considerably.

Each clip was then processed through a MetaHuman Performance asset \cite{epic_metahuman_performance} by assigning the imported capture data and a target MetaHuman. After processing, the resulting facial animation could be exported and used directly inside Sequencer. This made it possible to build the demo around a longer concatenated sequence of short emotional performances rather than leaving each clip as a separate technical test.

\begin{figure}[H]
    \centering
    \includegraphics[width=0.82\linewidth]{Demo3.png}
    \caption{Facial performance transfer into the main MetaHuman with pose correction.}
\end{figure}

\subsection{Application to Multiple MetaHumans}

Applying the processed facial performance to more than one MetaHuman allowed me to compare how the same source expression reads across different facial structures and appearances. Although the underlying solve is the same in principle, the perceived result changes depending on the target face shape and proportions.

For subtle or moderate expressions, the transfer was generally convincing. For stronger expressions, some nuance could be lost or slightly exaggerated, and the quality became more dependent on the target character and the source video.

Looking at the source clips directly against the generated facial animation made that even clearer. Core facial behaviour such as mouth articulation, eye motion, and broad expression timing transferred consistently, and because all clips came from the same actor speaking the same line, the differences seen in the output could be read mainly as differences in the target MetaHuman rather than inconsistencies in the source data. The timing of the animation stayed stable across both characters, but the way the emotion was perceived still shifted depending on the face it was applied to. Anger and surprise, in particular, could feel slightly stronger or harsher on some face shapes, while subtler expressions lost a little fine detail. Lip motion and phoneme timing were generally good, but still not exact enough to completely disappear as a solved process. That was probably the most useful comparison in the whole demo, because it showed that consistent transfer is not the same thing as identical emotional readability.

\subsection{Critical Evaluation}

Compared with a more manual facial pipeline, the combination of Live Link import and MetaHuman Performance processing gives a fairly direct path from source video to usable animation \cite{epic_live_link,epic_metahuman_performance}. That made it especially useful for quick experimentation and short emotional clips, where the aim was to compare performances rather than spend most of the time building setup around them.

At the same time, this is also the demo with the least direct editing control. Once the performance is solved, there is far less room to shape the result in the same hands-on way as in the body-animation demos, so the quality depends heavily on both the source video and the target face. That became especially obvious in the comparison between the two MetaHumans used here. My own MetaHuman has a chubbier face with fuller cheeks and lips, and the eyes are also more Asian in shape, so even though the solve remained usable, some expressions were softened or reduced when transferred and ended up feeling more adapted to my facial structure than directly identical to the source. I would place that translation at roughly 75\% overall. The alternative MetaHuman, by contrast, was a female character with a much sleeker face, and the expressions generally read more clearly on her despite the source clips coming from a male actor. On that face the translation felt stronger overall, and I would rate it closer to 90\%.

Because of that, this demo feels less like manual motion editing and more like evaluation of the pipeline itself. Its value comes from selecting suitable source clips, processing them efficiently, sequencing them clearly, and then comparing the output across different target characters. The direct comparison between the source video, my self-avatar, and the alternative MetaHuman is the strongest part here.
